
\clearpage
\chapter{אימון ואופטימיזציה}

\section{נתוני אימון סינתטיים}

לאימון המודל נוצרו נתונים סינתטיים: עבור כל דגימת אימון נבחרו באקראי \(K \in \{1,2,3\}\) מרכיבים עם משרעות \(A_k \sim U(0.1, 1.0)\), תדרים \(f_k \sim U(10, 500)\) הרץ ופאזות \(\phi_k \sim U(0, 2\pi)\). רעש גאוסי נוסף ברמות \textenglish{SNR} שבין \(-5\) ל-\textenglish{20 dB}.

\section{פונקציית אובדן ואסטרטגיית אימון}

פונקציית האובדן משלבת \textenglish{SI-SNR} ו-\textenglish{L1} על פרמטרים:

\begin{equation}
\mathcal{L} = -\text{SI-SNR}(s, \hat{s}) + \lambda \sum_{k} (|A_k - \hat{A}_k| + |f_k - \hat{f}_k|)
\end{equation}

עם \(\lambda = 0.01\). האופטימיזציה בוצעה באמצעות \textenglish{AdamW} עם קצב למידה \textenglish{$3 \times 10^{-4}$} ותזמון \textenglish{cosine annealing}. הושתמש בטכניקת \textenglish{curriculum learning}: בשלב הראשון נבחרו אותות עם \textenglish{SNR} גבוה (\textenglish{$> 15$ dB}), ובשלב השני הורחב לכלל ה-\textenglish{SNR} בטווח המוגדר.

\section{רגולריזציה ומניעת התאמת יתר}

\textenglish{Dropout} בשיעור \textenglish{0.2} הוחל על שכבות ה-\textenglish{LSTM} \cite{park2024hybrid}. בנוסף, \textenglish{weight decay} של \textenglish{$10^{-5}$} ו-\textenglish{gradient clipping} עם סף \textenglish{1.0} שיפרו את יציבות האימון. ניסויי \textenglish{ablation} הראו כי \textenglish{curriculum learning} מוסיף כ-\textenglish{0.8 dB} ל-\textenglish{SI-SNR} הסופי.
